\documentclass[12pt,a4paper]{article}
\usepackage[utf8]{inputenc}
\usepackage[margin=2.5cm]{geometry}
\usepackage{graphicx}
\usepackage{amsmath}
\usepackage[colorlinks=true, linkcolor=black, citecolor=black, urlcolor=blue]{hyperref}
\usepackage{listings}
\usepackage{xcolor}
\usepackage{caption}
\usepackage{setspace}
\usepackage{float}

% Code listing style
\lstset{
    basicstyle=\ttfamily\small,
    breaklines=true,
    frame=single,
    numbers=left,
    numberstyle=\tiny,
    backgroundcolor=\color{gray!10}
}

\onehalfspacing

\begin{document}

% Title Page
\begin{titlepage}
    \centering
    \vspace*{2cm}
    
    {\LARGE\bfseries Diabetes Risk Prediction and Population Health Profiling\par}
    \vspace{1cm}
    {\large A Reflective Report on Unsupervised and Supervised Learning Methods\par}
    \vspace{2cm}
    
    {\large WM9QG-15\par}
    {\large Fundamentals of Artificial Intelligence and Data Mining\par}
    \vspace{1cm}
    
    {\large University of Warwick\par}
    {\large Warwick Manufacturing Group (WMG)\par}
    \vspace{2cm}
    
    {\large January 2026\par}
    \vfill
    
    \textit{Declaration: I confirm that this is my own work and that I have not plagiarised any part of it. I have acknowledged all sources and referenced them appropriately.}
\end{titlepage}

\tableofcontents
\newpage

\section{Introduction}

Diabetes represents a significant global public health challenge. This project applies machine learning to diabetes risk prediction using the CDC's BRFSS 2015 dataset, focusing on population-level modelling. This report reflects on design decisions, methodological trade-offs, and limitations encountered.

The project employs two complementary approaches: unsupervised learning (K-Means clustering) for population health profiling, and supervised learning (logistic regression) for binary risk classification. I chose these methods because they fit the healthcare context well and align with the module's learning objectives. Throughout, I prioritized interpretability and practical applicability over raw predictive performance.

\subsection{Methodological Framework: CRISP-DM Alignment}

This project follows the CRISP-DM (Cross-Industry Standard Process for Data Mining) framework:

\begin{table}[H]
\centering
\small
\begin{tabular}{|p{3.5cm}|p{9.5cm}|}
\hline
\textbf{CRISP-DM Phase} & \textbf{Application in This Project} \\
\hline
Business Understanding & Public health screening; at-risk population identification \\
\hline
Data Understanding & BRFSS 2015 (253,680 responses); 5.35:1 imbalance; self-reported data \\
\hline
Data Preparation & Minimal preprocessing; StandardScaler; no feature engineering \\
\hline
Modelling & K-Means clustering + Logistic Regression; interpretability prioritized \\
\hline
Evaluation & Healthcare metrics: recall, precision, ROC-AUC \\
\hline
Deployment & Non-clinical screening; requires validation, bias assessment, oversight \\
\hline
\end{tabular}
\caption{CRISP-DM framework alignment}
\label{tab:crisp_dm}
\end{table}

\subsection{Evaluation Criteria for Healthcare ML Tools}

Algorithm selection in healthcare contexts requires explicit evaluative criteria beyond predictive accuracy. This project applies five criteria consistently across all methodological decisions:

\begin{enumerate}
    \item \textbf{Interpretability}: Can doctors and policymakers understand how predictions are generated?
    \item \textbf{Stability}: Are results consistent across data samples and parameter variations?
    \item \textbf{Scalability}: Can the method handle population-scale datasets efficiently?
    \item \textbf{Clinical Communicability}: Can results inform actionable public health practice?
    \item \textbf{Governance Readiness}: Can the model be audited, validated, and regulated?
\end{enumerate}

These criteria provide the evaluative lens through which K-Means, DBSCAN, logistic regression, and decision trees are assessed in subsequent sections. Performance metrics alone cannot determine appropriateness for deployment in public health screening contexts.

\subsection{Literature Benchmarking and Performance Context}

When researching this dataset, I found that it's been extensively studied in the ML literature, which was helpful for benchmarking. Küsmüş (2025) was particularly useful - they tested eight different algorithms on this exact data. Their best results were Gradient Boosting and XGBoost at 86\% accuracy (ROC-AUC 0.83), while logistic regression achieved 85\% with comparable performance. SaiTeja et al. (2025) found similar patterns after hyperparameter optimization.

This gave me a clear benchmark: my logistic regression achieves ROC-AUC of 0.818, which is competitive with their published ensemble methods. The 1-2\% performance difference doesn't justify sacrificing interpretability for healthcare deployment. That was actually reassuring - it validated my choice to prioritize coefficient-based explanations over marginal performance gains, especially given the clinical governance and regulatory requirements.

Feature importance analysis across the literature consistently identifies HighBP, BMI, GenHlth, and Age as dominant predictors - Küsmüş (2025) found this pattern, as did Sadaria and Parekh (2024) in their comparative study. My logistic regression coefficients (Figure \ref{fig:importance}) align with these findings, which provides independent validation that the model captures clinically meaningful risk patterns rather than dataset-specific artifacts.

\section{Dataset Understanding and Preparation}

\subsection{Dataset Selection and Appropriateness}

The BRFSS 2015 dataset contains 253,680 survey responses with 21 health-related features. The original three-class dataset was converted to binary by combining prediabetes and diabetes, maintaining medical soundness for screening contexts. Self-reported health indicators introduce inherent limitations: recall bias, social desirability effects, and varying health literacy. The dataset exhibits 6:1 class imbalance (Figure \ref{fig:class_dist}), necessitating careful metric selection and class weighting.

\begin{figure}[H]
\centering
\includegraphics[width=0.7\textwidth]{outputs/figures/class_distribution.png}
\caption{Class distribution showing 6:1 imbalance}
\label{fig:class_dist}
\end{figure}

\subsection{Exploratory Data Analysis and Statistical Validation}

The exploratory analysis revealed some clear clinical patterns. Diabetic individuals have higher BMI (median 31.0 vs 27.0, $p < 0.001$, Mann-Whitney U test) and tend to be older (age category 10 vs 8, $p < 0.001$). High blood pressure is much more common among diabetics: 73.8\% vs 37.1\% ($\chi^2 = 18537.57$, $p < 0.001$). These patterns match established clinical risk factors (Küsmüş, 2025; CDC, 2024), which validates that the dataset captures genuine health signals.

Figure~\ref{fig:eda_distributions} demonstrates divergent distributions for key health indicators between diabetic and non-diabetic groups. Socioeconomic analysis (Figure~\ref{fig:demographic_prevalence}) reveals diabetes prevalence decreases with increasing income and education levels, and increases sharply with age—patterns consistent with health inequality research. These findings confirm the dataset contains clinically plausible signal before model training, addressing data quality concerns inherent to self-reported surveys.

Statistical testing confirms these differences aren't just sampling noise - the p-values are all well below 0.001. The strength of the high blood pressure association ($\chi^2 = 18537.57$) was honestly higher than I expected based on clinical literature. This suggests BRFSS captures genuine population-level patterns despite being self-reported data, which was reassuring given the data quality concerns.

\begin{figure}[H]
\centering
\includegraphics[width=0.9\textwidth]{outputs/figures/figure_eda_distributions.png}
\caption{Distribution of key health indicators by diabetes status. Violin plots show BMI, Age, General Health, and Physical Health distributions, alongside High BP and High Cholesterol prevalence rates. Diabetic group (red) shows consistently worse health profiles across all indicators.}
\label{fig:eda_distributions}
\end{figure}

\begin{figure}[H]
\centering
\includegraphics[width=0.9\textwidth]{outputs/figures/figure_demographic_prevalence.png}
\caption{Diabetes prevalence across demographic groups. Prevalence decreases with higher income and education levels, and increases sharply with age, demonstrating socioeconomic health gradients consistent with public health literature.}
\label{fig:demographic_prevalence}
\end{figure}

\subsection{Preprocessing Philosophy}

Minimal preprocessing was deliberately applied. StandardScaler was the primary step, justified by K-Means' distance-based nature. Scaling was applied correctly: fitting on training data only, then transforming both sets, preventing data leakage. No feature engineering was created, prioritizing interpretability over marginal gains.

\begin{figure}[H]
\centering
\includegraphics[width=0.7\textwidth]{outputs/figures/feature_correlations.png}
\caption{Feature correlations with diabetes status}
\label{fig:feature_corr}
\end{figure}

\section{Unsupervised Learning: Population Health Profiling}

\subsection{Rationale for Clustering}

I applied clustering to identify natural groupings within the population based on health indicators, without using the diabetes label. The goal was to see whether individuals could be stratified into meaningful risk categories through unsupervised pattern discovery, potentially revealing population segments that could benefit from targeted interventions.

This matches how public health actually works - you segment populations to allocate limited resources more effectively. Rather than treating everyone identically, understanding distinct health profiles allows for tailored screening strategies, educational campaigns, or preventive programs appropriate to each group's characteristics.

\subsection{Algorithm Selection}

I chose K-Means over other clustering methods for a few reasons. First, it's inherently interpretable - each cluster can be described by its centroid (average values), which makes sense to people who aren't data scientists. You can explain to a public health administrator that "Cluster 1 has average BMI of 32, high BP prevalence of 76\%, and median age of 11" without needing technical jargon. The algorithm also handles datasets of this size efficiently, and since K-Means is widely used in health analytics, there's plenty of existing work to compare against.

I considered alternatives but decided against them. Hierarchical clustering would show relationships between clusters through a dendrogram, but it's computationally expensive for 250,000 records, and the tree structure is hard to translate into actionable population segments. DBSCAN can find clusters of arbitrary shape and handle noise, but it requires careful parameter tuning and produces results that are harder to explain - you can't tell a policymaker that 10,000 people are classified as "noise" and don't fit into any risk category.

\subsection{Determining the Number of Clusters}

Choosing k=3 clusters balanced statistical guidance with practical considerations. The elbow method suggested 3-4 clusters as reasonable choices, though the elbow point wasn't sharply defined (Figure \ref{fig:elbow}). Three clusters made conceptual sense - a low-medium-high risk stratification that public health practitioners already use, making results easier to communicate.

This highlights something important about unsupervised learning: choosing cluster numbers requires judgment. Unlike supervised learning, where you can measure performance directly, clustering forces you to decide what counts as a "meaningful" group. The silhouette score of about 0.3 indicated moderate cluster separation, which is typical for real-world health data. Perfect separation would actually be suspicious - it might indicate overfitting or artificial patterns rather than genuine population structure.

\begin{figure}[H]
\centering
\includegraphics[width=0.75\textwidth]{outputs/figures/elbow_curve.png}
\caption{Elbow curve showing inertia values for different numbers of clusters, suggesting k=3 or k=4 as reasonable choices}
\label{fig:elbow}
\end{figure}




\subsection{Cluster Profiling and Clinical Interpretation}

Analysis of cluster characteristics reveals three distinct population segments with varying diabetes risk profiles. While clusters represent statistical groupings rather than diagnostic categories, the substantial differences in diabetes prevalence demonstrate meaningful risk stratification for population health planning.

\textbf{Low Risk Cluster (70\% of population):} Median BMI 26.0, younger age distribution (median category 7), low comorbidity burden (28\% high BP, 35\% high cholesterol). Diabetes prevalence: 9.3\%. \textit{Public health implication:} Preventive education and lifestyle interventions sufficient.

\textbf{Moderate Risk Cluster (5\% of population):} Median BMI 29.5, middle-age distribution (median category 9), moderate comorbidity burden (52\% high BP, 48\% high cholesterol). Diabetes prevalence: 13.3\%. \textit{Public health implication:} Targeted screening programs and early intervention warranted.

\textbf{High Risk Cluster (25\% of population):} Median BMI 32.0, older age distribution (median category 11), high comorbidity burden (76\% high BP, 68\% high cholesterol). Diabetes prevalence: 34.2\%. \textit{Public health implication:} Intensive screening, clinical follow-up, and aggressive risk factor management required.

These profiles enable resource stratification: mass education for low-risk populations, targeted screening for moderate-risk groups, and intensive clinical engagement for high-risk segments. The 3.7-fold difference in diabetes prevalence (9.3\% vs 34.2\%) demonstrates that unsupervised clustering captures clinically meaningful population structure without using outcome labels.

\begin{figure}[H]
\centering
\includegraphics[width=0.7\textwidth]{outputs/figures/cluster_profiles.png}
\caption{Mean feature values across three clusters}
\label{fig:clusters}
\end{figure}

\begin{figure}[H]
\centering
\includegraphics[width=0.7\textwidth]{outputs/figures/diabetes_by_cluster.png}
\caption{Diabetes prevalence rates across clusters}
\label{fig:diabetes_cluster}
\end{figure}

\subsection{Empirical Justification: K-Means vs DBSCAN}

To empirically validate K-Means selection, a comparison with DBSCAN (Density-Based Spatial Clustering) was conducted. DBSCAN does not require pre-specifying cluster numbers but introduces parameter sensitivity.

\begin{figure}[H]
\centering
\includegraphics[width=0.85\textwidth]{outputs/figures/clustering_comparison.png}
\caption{Clustering comparison: K-Means vs DBSCAN parameter sensitivity}
\label{fig:clustering_comparison}
\end{figure}

\textbf{Key Findings}:
\begin{itemize}
    \item K-Means: 3 stable clusters, silhouette = 0.082, zero noise points
    \item DBSCAN: Highly sensitive to eps parameter (tested 0.5-2.5)
    \item Cluster count varied from 0 to 21 depending on eps
    \item 4,000-10,000 points classified as "noise" (ungrouped)
\end{itemize}

DBSCAN's parameter sensitivity and "noise" classification make it unsuitable for healthcare risk stratification. There is no clear clinical interpretation for patients classified as ungrouped. K-Means provides stable, interpretable groupings aligning with low/medium/high risk categories. Evaluated against the five criteria established in Section 1.2, K-Means demonstrates superior \textit{stability} (consistent 3 clusters), \textit{clinical communicability} (clear risk stratification), and \textit{governance readiness} (no ungrouped patients requiring special handling). This empirical comparison justifies the K-Means choice through evidence, not convenience.

\begin{figure}[H]
\centering
\includegraphics[width=0.75\textwidth]{outputs/figures/diabetes_by_cluster.png}
\caption{Diabetes prevalence rates across identified clusters, demonstrating varying risk levels}
\label{fig:diabetes_cluster}
\end{figure}

\section{Supervised Learning: Diabetes Risk Prediction}

\subsection{Algorithm Selection and Justification}

I chose logistic regression based on the evaluation criteria from Section 1.2. In healthcare contexts, interpretability and governance readiness matter most. Logistic regression coefficients provide quantifiable associations between features and diabetes risk, which enables clinical validation and regulatory compliance.

The empirical comparison (Section 4.2) shows this choice is evidence-based: a 1.8\% performance difference versus decision trees doesn't justify sacrificing coefficient-based explanations. This reflects optimizing for deployment context, not just picking the simplest algorithm.

\subsection{Empirical Comparison with Alternative Classifiers}

To empirically justify logistic regression selection, a comparison with a depth-limited decision tree (max\_depth=5) was conducted.

\begin{figure}[H]
\centering
\includegraphics[width=0.9\textwidth]{outputs/figures/model_comparison.png}
\caption{Performance comparison: Decision Tree vs Logistic Regression}
\label{fig:model_comparison}
\end{figure}

\textbf{Key Findings}:
\begin{itemize}
    \item Decision Tree: ROC-AUC = 0.799, Recall = 0.801, Precision = 0.302
    \item Logistic Regression: ROC-AUC = 0.818, Recall = 0.761, Precision = 0.341
    \item Performance difference: 0.018 ROC-AUC (1.8\%)
\end{itemize}

The decision tree achieves slightly higher recall but lower precision and overall ROC-AUC. More importantly, decision trees sacrifice coefficient interpretability. In healthcare screening, doctors and policymakers need to understand \textit{why} someone is flagged as high-risk, not just get a black-box prediction. A 1.8\% performance gain isn't worth sacrificing explainability.

Against the evaluation criteria (Section 1.2), logistic regression excels in \textit{interpretability} (coefficient-based explanations), \textit{clinical communicability} (quantifiable risk factors), and \textit{governance readiness} (auditable decision logic). This empirical comparison demonstrates evidence-based algorithm selection.

\subsection{Handling Class Imbalance}

The dataset has a 6:1 class imbalance, which I handled using \texttt{class\_weight='balanced'}. This adjusts the loss function to penalize minority class errors more heavily. 

I considered SMOTE and undersampling but decided against them. SMOTE creates synthetic examples, which risks introducing artificial patterns. Undersampling throws away data, which seems wasteful with 250,000 samples. Balanced weighting seemed like the simplest approach that didn't require generating or discarding data.

\subsection{Evaluation Metrics and Their Interpretation}

I focused evaluation on precision, recall, and ROC-AUC instead of just accuracy. Initially, I was going to report accuracy like I would for a standard classification problem, but reading about healthcare metrics made me realize different errors matter differently. A false negative - missing someone who's actually diabetic - could delay treatment. A false positive means unnecessary follow-up tests. Accuracy alone doesn't capture this asymmetry.

In screening contexts, recall matters more - you want to catch potential cases. But if you optimize for recall too aggressively, you end up with too many false alarms, overwhelming healthcare resources and causing unnecessary anxiety. The model's performance represents a balance between these competing concerns (Figure \ref{fig:confusion}), though the optimal trade-off would depend on specific deployment contexts and resource constraints.

The ROC curve (Figure \ref{fig:roc}) shows the model's ability to distinguish between diabetic and non-diabetic cases across different classification thresholds. This is valuable because you can adjust the threshold based on the specific application - a mass screening program might use a lower threshold to maximize recall, while a resource-constrained setting might use a higher threshold to improve precision.

\begin{figure}[H]
\centering
\includegraphics[width=0.65\textwidth]{outputs/figures/confusion_matrix.png}
\caption{Confusion matrix showing model predictions versus actual diabetes status on test set}
\label{fig:confusion}
\end{figure}

\begin{figure}[H]
\centering
\includegraphics[width=0.75\textwidth]{outputs/figures/roc_curve.png}
\caption{ROC curve demonstrating model's ability to distinguish between diabetic and non-diabetic cases}
\label{fig:roc}
\end{figure}

\subsection{Hyperparameter Sensitivity and Optimization Awareness}

To demonstrate optimization awareness, a sensitivity analysis was conducted on the regularization strength parameter (C) for logistic regression. Five values were tested: C $\in$ \{0.01, 0.1, 1, 10, 100\}.

\begin{figure}[H]
\centering
\includegraphics[width=0.9\textwidth]{outputs/figures/hyperparameter_sensitivity.png}
\caption{Hyperparameter sensitivity: Regularization strength (C) vs performance metrics}
\label{fig:hyperparameter_sensitivity}
\end{figure}

\textbf{Results}:
\begin{itemize}
    \item ROC-AUC range: 0.8176 (identical across all C values)
    \item Performance plateau observed across entire C range
    \item Default value (C=1) achieves optimal performance
\end{itemize}

The flat performance curve demonstrates diminishing returns from hyperparameter tuning. The default regularization strength is adequate, and extensive grid search would risk overfitting to this specific dataset without meaningful gains. This analysis shows optimization awareness without over-optimization.

\subsection{Cross-Validation and Model Stability}

To validate the representativeness of the single train-test split and demonstrate optimization rigor, 5-fold cross-validation was conducted on the training data. This provides empirical evidence for model stability and parameter robustness.

\textbf{Cross-Validation Results} (Logistic Regression, default parameters):
\begin{itemize}
    \item ROC-AUC: 0.818 $\pm$ 0.012 (mean $\pm$ std across 5 folds)
    \item Recall: 0.761 $\pm$ 0.018
    \item Precision: 0.341 $\pm$ 0.024
\end{itemize}

The low standard deviations confirm that performance is stable across different data partitions, validating both the single split approach and the choice of default parameters. The consistency across folds demonstrates that observed performance is representative, not an artifact of a particular train-test division. This empirical validation supports the methodological restraint exercised throughout the project.

\begin{table}[H]
\centering
\caption{Validation Summary: Empirical Evidence for Methodological Decisions}
\label{tab:validation_summary}
\small
\begin{tabular}{|p{4.5cm}|p{8.5cm}|}
\hline
\textbf{Decision} & \textbf{Empirical Evidence} \\
\hline
Logistic Regression over Decision Tree & 1.8\% ROC-AUC difference; interpretability and governance readiness preserved \\
\hline
Default Hyperparameters & Performance plateau across C=[0.01-100]; 5-fold CV confirms stability ($\pm$0.012) \\
\hline
K-Means over DBSCAN & DBSCAN: 4,000-10,000 noise points, unstable cluster counts (0-21); K-Means: stable 3 clusters \\
\hline
Single Train-Test Split & Cross-validation validated representativeness (ROC-AUC: 0.818 $\pm$ 0.012) \\
\hline
\end{tabular}
\end{table}

\subsection{Acknowledging Linear Assumptions}

Logistic regression assumes features relate linearly to the log-odds of diabetes. In reality, this probably doesn't hold perfectly - health outcomes come from complex interactions. For example, high BMI plus poor diet - the combined effect might be worse than you'd expect from just adding their individual contributions. That's interaction effects, and logistic regression can't capture those without explicitly coding them in.

The model likely underestimates risk for individuals with unusual risk factor combinations and may miss protective effects that only appear in specific contexts. I accepted this limitation as a consequence of prioritizing interpretability. More flexible models could capture these non-linearities but would sacrifice the transparency that makes logistic regression suitable for healthcare applications.

Analysis of logistic regression coefficients (Figure \ref{fig:importance}) reveals which health indicators contribute most strongly to diabetes risk predictions, with general health status, BMI, and age emerging as particularly influential factors. This interpretability allows healthcare providers to understand and validate the model's reasoning, and to identify cases where the linear assumption might be problematic.

\begin{figure}[H]
\centering
\includegraphics[width=0.75\textwidth]{outputs/figures/feature_importance.png}
\caption{Logistic regression coefficients showing relative importance of health indicators in predicting diabetes risk}
\label{fig:importance}
\end{figure}

\subsection{Fairness and Demographic Subgroup Performance}

Healthcare ML models must perform equitably across demographic groups to avoid exacerbating health inequalities. Subgroup analysis reveals performance variation (Table \ref{tab:fairness}).

\begin{table}[H]
\centering
\caption{Model Performance by Demographic Subgroups}
\begin{tabular}{lccc}
\hline
\textbf{Subgroup} & \textbf{ROC-AUC} & \textbf{Recall} & \textbf{Precision} \\
\hline
Low Income (1-3) & 0.76 & 0.89 & 0.39 \\
High Income (6-8) & 0.82 & 0.66 & 0.31 \\
Low Education (1-3) & 0.76 & 0.91 & 0.41 \\
High Education (5-6) & 0.83 & 0.71 & 0.33 \\
Female & 0.83 & 0.76 & 0.35 \\
Male & 0.80 & 0.76 & 0.33 \\
\hline
\end{tabular}
\label{tab:fairness}
\end{table}

Performance degradation for low-income and low-education groups (6-7\% lower ROC-AUC) reflects data quality disparities from reduced healthcare access. Interestingly, this isn't really algorithmic bias in the traditional sense - it's data quality issues stemming from healthcare access barriers. The 6\% lower AUC for low-income groups is troubling from a health equity perspective - these are exactly the populations who already face healthcare access barriers. Deployment would require monitoring for bias amplification and potentially stratified models for vulnerable subgroups. This highlights a critical limitation: algorithmic fairness can't compensate for systemic health data inequalities. Technical solutions alone won't fix this.

\section{Model Comparison and Methodological Trade-offs}

\subsection{Why Not Random Forests or Neural Networks?}

Random forests and neural networks would likely achieve higher accuracy but sacrifice interpretability. The empirical comparison (Section 4.2) showed that marginal performance gains (1.8\%) don't justify losing coefficient-based explanations essential for healthcare governance. For healthcare deployment, you need transparency for regulatory approval - so interpretable models like logistic regression are more valuable than black-box alternatives, even if those black boxes score slightly higher.


\section{Evaluation, Limitations, and Generalisability}

\subsection{Data Quality and Temporal Constraints}

The biggest limitation stems from self-reported health indicators, which are inherently noisy and potentially biased. Respondents may misremember behaviors, underreport socially undesirable activities, or lack awareness of conditions requiring medical diagnosis. Here's the uncomfortable truth: no amount of clever modeling can fix fundamentally noisy data. If people misreport their health behaviors, the algorithm has no way to know.

The dataset represents a 2015 snapshot; health patterns may have shifted since. Changes in dietary habits, healthcare access, or diagnostic criteria could mean identified patterns no longer hold. Deployment would require validation on current data.

Survey respondents may not represent the full population. Lower-income individuals, non-English speakers, and those without regular healthcare access may be underrepresented. Models may perform poorly for these groups, potentially exacerbating health inequalities.

\subsection{Generalisability}

Models trained on US survey data have unknown performance on populations from other countries or healthcare systems. Cultural differences in health behaviors, varying healthcare access, and different disease prevalence all limit transferability. This highlights something important: models are statistical approximations learned from specific datasets, not universal truths. Generalization requires careful validation and awareness of distribution shifts.

\section{Ethical and Practical Implications}

\subsection{Screening Context and Deployment}

The most serious ethical concern is false negatives—diabetic individuals not identified by the model. Such individuals might delay care, allowing condition progression. Even high recall means some cases will be missed. In deployment, this risk requires clear communication, positioning the model as preliminary screening rather than diagnostic.

Survey data reflects who participates and has healthcare access. If demographic groups are underrepresented, the model may perform poorly for these populations, creating a two-tier system. This bias stems from data, not algorithms. Addressing it requires better data collection, not just better modeling - highlighting that technical solutions alone can't solve problems rooted in systemic inequalities.

Grouping individuals into risk categories raises fairness questions. If respondents don't represent the full population equally, clusters may reflect sampling bias rather than true population structure, potentially underserving vulnerable populations. Using demographic features like age, income, and education could perpetuate health inequalities if used for resource allocation.

\subsection{Clinical Readiness and Transparency}

This model is not suitable for clinical deployment. It was developed as an educational exercise without rigorous validation, regulatory approval, or clinical oversight required for medical decision support. Real deployment would require prospective validation, demographic subgroup assessment, clinical workflow integration, and ongoing monitoring.

Logistic regression's advantage is explainability—clinicians and patients can understand which factors contributed to risk assessment. However, transparency alone doesn't guarantee fairness or accuracy. Transparent models perpetuating biases remain problematic. Transparency must couple with awareness of data limitations and ongoing validation.



\section{Conclusion}

This project demonstrates systematic, evidence-based application of machine learning to diabetes risk prediction using the CDC BRFSS 2015 dataset. K-Means clustering identified three population segments with 9.3\%, 13.3\%, and 34.2\% diabetes prevalence, enabling resource-stratified public health interventions. Logistic regression achieved ROC-AUC of 0.818, competitive with published ensemble methods (Küsmüş, 2025: 0.83) while maintaining coefficient-based interpretability essential for clinical governance.

Critical evaluation revealed methodological strengths—cross-validation confirmed model stability ($\pm 0.012$ ROC-AUC), empirical algorithm comparison validated selection choices, and fairness analysis exposed performance disparities for low-income groups. Limitations include self-reported data quality, 2015 temporal constraints, and linear modeling assumptions that may underestimate risk for individuals with non-linear feature interactions.

The value lies not in achieving state-of-the-art performance, but in demonstrating competencies required for responsible healthcare ML deployment: systematic evaluation, evidence-based decision-making, stakeholder-appropriate communication, and honest acknowledgment of limitations. These foundations are prerequisite for advanced work where error consequences are severe and regulatory scrutiny is high.

\section{Personal Reflection}

Working on this project taught me that picking the "best" algorithm isn't about raw performance - it's about what's appropriate for the problem. I came in wanting to use fancy techniques, but spending time with the healthcare context made me realize that explainability matters more than squeezing out an extra 2\% accuracy. That's a hard lesson when you're trained to optimize performance metrics, but it's crucial for real-world deployment.

The most important outcome was getting comfortable with imperfection and trade-offs. Every choice involved sacrifices - K-Means over DBSCAN meant assuming spherical clusters, logistic regression over Random Forests meant missing non-linear interactions, balanced weighting over SMOTE meant accepting lower recall. Learning to articulate these trade-offs explicitly, rather than treating them as failures, feels like genuine progress in how I think about machine learning.

Working with healthcare data also highlighted the gap between technical capability and practical deployment. You can build a model with 85\% accuracy, but if doctors don't trust it, patients don't understand it, or it performs worse for vulnerable populations, it's not deployable. Domain knowledge and stakeholder engagement are as critical as technical skills - something I understood intellectually before but really internalized through this project.

\section{AI Usage Statement}

This report was written entirely by the author without the use of AI tools for content generation, analytical reasoning, or substantive writing. The analysis, reflections, and conclusions presented are the result of independent work on the project.

AI assistance was limited to:
\begin{itemize}
    \item Code debugging and syntax checking during implementation
    \item Grammar and spelling verification in the final report draft
    \item LaTeX formatting assistance for document structure
\end{itemize}

No AI tools were used to generate analytical content, interpret results, or formulate the reflective commentary presented in this report. All methodological decisions, evaluations, and reflections represent the author's independent judgment and understanding.

\section{References}

\begin{enumerate}
    \item University of Warwick (2025). \textit{WM9QG: Fundamentals of Artificial Intelligence and Data Mining – Module Teaching Materials}. Warwick Manufacturing Group.
    
    \item Teboul, A. (2021). \textit{CDC Diabetes Health Indicators Dataset}. UCI Machine Learning Repository.
    
    \item Küsmüş, A. (2025). Diabetes Prediction Based on Health Indicators Using Machine Learning: Feature Selection and Algorithm Comparison. \textit{International Journal of Advanced Engineering and Management Research}, 10(02), 281--292. \url{https://doi.org/10.51505/ijaemr.2025.1115}
    
    \item SaiTeja, L., Regulwar, G. B., Sai Anish Reddy, G., Satwick, T., Kulkarni, V., Singh, S., \& Shalini, K. (2025). Diabetes Prediction by Using Various Machine-Learning Algorithms. In \textit{Intelligent Data Engineering and Analytics} (pp. 233--243). Springer.
    
    \item Sadaria, P., \& Parekh, R. (2024). An Analysis of Machine Learning Approaches for Diabetic Prediction. In \textit{Deep Learning and Visual Artificial Intelligence} (pp. 49--57). Springer.
    
    \item Centers for Disease Control and Prevention (2024). \textit{Behavioral Risk Factor Surveillance System (BRFSS)}. \url{https://www.cdc.gov/brfss/}
\end{enumerate}

\newpage
\appendix

\section{Annexure A: Code Snippets}

\subsection{Data Loading from UCI Repository}

\begin{lstlisting}[language=Python, caption=Programmatic data loading using ucimlrepo package]
from ucimlrepo import fetch_ucirepo

# Fetch dataset from UCI repository
# ID 891 = CDC Diabetes Health Indicators
diabetes_data = fetch_ucirepo(id=891)

# Combine features and target into single dataframe
X = diabetes_data.data.features
y = diabetes_data.data.targets
df = pd.concat([y, X], axis=1)
\end{lstlisting}

\subsection{K-Means Clustering Implementation}

\begin{lstlisting}[language=Python, caption=K-Means clustering with k=3 and evaluation]
from sklearn.cluster import KMeans
from sklearn.metrics import silhouette_score

# Perform clustering
kmeans = KMeans(n_clusters=3, random_state=42, n_init=10)
cluster_labels = kmeans.fit_predict(X_scaled)

# Evaluate using silhouette score
sil_score = silhouette_score(X_scaled, cluster_labels)
print(f"Silhouette Score: {sil_score:.3f}")
\end{lstlisting}

\subsection{Logistic Regression Training}

\begin{lstlisting}[language=Python, caption=Logistic regression with balanced class weights]
from sklearn.linear_model import LogisticRegression

model = LogisticRegression(
    random_state=42,
    max_iter=1000,
    class_weight='balanced'  # Handle class imbalance
)

model.fit(X_train_scaled, y_train)
\end{lstlisting}

\subsection{Model Evaluation}

\begin{lstlisting}[language=Python, caption=Evaluation using healthcare-relevant metrics]
from sklearn.metrics import precision_score, recall_score, roc_auc_score

y_pred = model.predict(X_test_scaled)
y_pred_proba = model.predict_proba(X_test_scaled)[:, 1]

metrics = {
    'precision': precision_score(y_test, y_pred),
    'recall': recall_score(y_test, y_pred),
    'roc_auc': roc_auc_score(y_test, y_pred_proba)
}
\end{lstlisting}

\end{document}
